\documentclass[11pt,a4paper]{article}

\usepackage[margin=2.5cm]{geometry}
\usepackage{hyperref}
\usepackage{amsmath,amssymb}
\usepackage{parskip}
\usepackage{microtype}
\usepackage{titlesec}

\titleformat{\section}{\large\bfseries}{\thesection}{0.6em}{}
\titleformat{\subsection}{\normalsize\bfseries}{\thesubsection}{0.6em}{}

\title{\textbf{Methodology: Branched Multi-Turn GRPO for Code-Policy Generation}}
\author{}
\date{}

\begin{document}
\maketitle

\section{Setup and notation}

The agent is a language model $\pi_\theta$ that, conditioned on a textual
prompt $s$ describing a control task, emits a Python policy. The base is
\texttt{Qwen2.5-Coder-3B-Instruct} with frozen weights $W$; $\theta$
parameterizes a LoRA adapter giving effective layer weights
$W' = W + AB^\top$. Interaction with the environment is mediated by a
chat template, so a trajectory is a sequence of role-typed messages.

A task $\mathcal{T}$ specifies environment source code, $S$ scenarios
$\{\sigma_1,\dots,\sigma_S\}$ that vary the seed and hidden parameters,
and a per-scenario reward. The scoring oracle runs a candidate policy
body in a sandboxed subprocess against each scenario and returns
\[
  R(\text{code};\mathcal{T}) \;=\; \sum_{i=1}^{S} r_i(\text{code};\sigma_i),
\]
with $r_i$ the cumulative episode reward (capped at \texttt{max\_steps}).
Validation failures, runtime errors, and timeouts on $\sigma_i$ score the
\emph{error floor} $r_i = -200$, calibrated below any running-but-failing
trajectory --- otherwise errored rollouts (which accumulate $0$) would
outrank running-but-poor rollouts that integrate negative distance
penalties, and the gradient would reward crashing.

\subsection{Multi-turn protocol ($N=2$)}

Each turn the model emits exactly one fenced code block tagged either
\texttt{policy} or \texttt{main}. A \texttt{policy} block defines
\texttt{def policy(state)}; the tool runs it across all scenarios and
returns the summed reward as a user-role message. A \texttt{main} block
is arbitrary scratchpad code with no environment access; the tool returns
truncated stdout. The final turn must be \texttt{policy}; \texttt{main}
or untagged final turns are rejected at the floor $-200\cdot S$.
Malformed intermediate turns are likewise floored and surfaced as a
textual rejection in the next user message.

The full trajectory is
\begin{equation}
  \tau \;=\; \big(s,\; y^{(1)},\; m,\; y^{(2)}\big),
  \label{eq:trajectory}
\end{equation}
where $y^{(1)}, y^{(2)}$ are the two completions and $m$ is the
user-role tool message returned after executing $y^{(1)}$. The terminal
reward $R(\tau) = R(\text{code in }y^{(2)};\mathcal{T})$ is the only
scalar reward; tool outputs are observations. The rationale for two
turns: $y^{(1)}$ either rehearses a draft \texttt{policy} (and reads
back its actual scenario reward) or runs a \texttt{main} scratchpad to
print intermediate quantities; $y^{(2)}$ is the scored
\texttt{policy} we deploy.

\section{Background: clipped surrogate without a critic}

For a single-turn rollout with completion $y$ of length $T$, the standard
clipped surrogate is
\begin{equation}
\mathcal{L}^{\text{PPO}}(\theta) \;=\;
  -\,\mathbb{E}_{(s,y)\sim q}\!\left[
    \frac{1}{T}\sum_{t=1}^{T}\,
    \min\!\big(\rho_t A_t,\;
               \mathrm{clip}(\rho_t, 1-\epsilon_{\text{lo}}, 1+\epsilon_{\text{hi}}) A_t\big)
  \right],
  \label{eq:ppo}
\end{equation}
with $\rho_t(\theta) = \pi_\theta(y_t \mid s, y_{<t})/q(y_t\mid s, y_{<t})$
and $q$ the (frozen) sampling distribution. PPO requires a learned
value network for $A_t$, awkward to fit alongside a LoRA-adapted policy
on a single GPU.

Group-relative variants~\cite{shao2024deepseekmath} drop the critic
in favor of a Monte Carlo per-prompt baseline: sample $G$ completions
$\{y^{(g)}\}$, score each, and broadcast a standardized reward to all
tokens of $y^{(g)}$,
\begin{equation}
A^{(g)}_t \;=\; \frac{R(y^{(g)}) - \bar R}{\widehat\sigma_R}
\quad\text{(constant in $t$).}
\end{equation}

We use the \emph{token-summed} variant of the clipped objective: drop
the per-row mean, sum the clipped term over all tokens in the global
batch, normalize once by the count of valid generated tokens,
\begin{equation}
\mathcal{L}(\theta) \;=\;
 -\,\frac{1}{N_{\text{tok}}}\sum_{b,t}
   \,\mathbf{1}\{\text{model token}\}\cdot
   \min\!\big(\rho^b_t A^b_t,\;\mathrm{clip}(\rho^b_t,\,1-\epsilon_{\text{lo}},\,1+\epsilon_{\text{hi}}) A^b_t\big),
\label{eq:loss}
\end{equation}
with $N_{\text{tok}} = \sum_{b,t}\mathbf{1}\{\text{model token in }(b,t)\}$.\footnote{This
matches the DAPO~\cite{yu2025dapo} and CISPO recipes; the per-row
averaging variants used by the original GRPO formulation, BNPO, and
Dr.GRPO instead divide by per-row token counts.}
The form is chosen for one structural reason: \eqref{eq:loss} is
\emph{additive over disjoint token subsets} --- for any partition of the
model-token positions, the loss equals the sum of partial losses with
the same $N_{\text{tok}}$. This makes the per-turn decomposition in
\S\ref{sec:two-pass} exact.

\section{Branched rollouts and per-turn credit assignment}
\label{sec:branched}

A direct extension of group-relative training to two-turn trajectories
draws $G$ samples per prompt, each containing both turns, and broadcasts
a single trajectory-level advantage to every model token in
$y^{(1)}\cup y^{(2)}$. This conflates turn-1 and turn-2 credit: a strong
turn-1 followed by a weak turn-2 sample receives the same penalty on
both turns.

\subsection{The $G_1 \times G_2$ tree}

We branch at the turn boundary, in the spirit of VinePPO's
intra-completion MC branching~\cite{kazemnejad2024vineppo}. For each
prompt $s$, sample $G_1$ turn-1 variants in one \texttt{model.generate}
call,
\[
  y^{(1)}_i \sim q(\cdot\mid s),\qquad i=1,\dots,G_1.
\]
For each $y^{(1)}_i$, execute the tool to obtain $m_i$, then sample
$G_2$ continuations in one \texttt{generate} call,
\[
  y^{(2)}_{i,j} \sim q(\cdot\mid s,\, y^{(1)}_i,\, m_i),\qquad j=1,\dots,G_2.
\]
Score each leaf to obtain rewards $\{R_{i,j}\}$. Total generate calls
per prompt: $1+G_1$, versus $2G_1 G_2$ for an unbranched $N=2$ rollout
producing the same number of leaves.

\subsection{Monte Carlo value estimates and per-turn advantages}

Two empirical values, both pure averages over the tree:
\begin{align}
  \widehat V_0          &\;=\; \frac{1}{G_1 G_2}\sum_{i,j} R_{i,j}, &
  \widehat V_1^{(i)}    &\;=\; \frac{1}{G_2}\sum_{j} R_{i,j}.
\end{align}
$\widehat V_0$ estimates the initial-state value (after the prompt,
before any model action); $\widehat V_1^{(i)}$ estimates the value
after committing to turn-1 variant $y^{(1)}_i$ and observing $m_i$.
With deterministic environments (given the seed) and $\gamma=1$, both
are unbiased MC estimates of the corresponding state values under $q$.

The per-turn advantages, broadcast to every token of the corresponding
turn, are
\begin{align}
  A^{(1)}_{i}      &\;=\; \widehat V_1^{(i)} \;-\; \widehat V_0,
  &\text{(over $y^{(1)}_i$)}, \label{eq:adv1}\\
  A^{(2)}_{i,j}    &\;=\; R_{i,j} \;-\; \widehat V_1^{(i)},
  &\text{(over $y^{(2)}_{i,j}$).} \label{eq:adv2}
\end{align}
The turn-2 advantage uses the \emph{conditional} mean given the turn-1
prefix rather than the global per-prompt mean: a poor turn-2 sample
drawn from a strong turn-1 is recognized as poor instead of inheriting
its prefix's above-average outcome.

Tool-message tokens are masked out via a separate \texttt{tool\_mask};
their advantage is set to zero. The per-token advantage vector for
trajectory $(i,j)$ is
\begin{equation}
  A_{i,j} \;=\; \big(\underbrace{A^{(1)}_{i},\dots,A^{(1)}_{i}}_{T_1},\;
                     \underbrace{0,\dots,0}_{T_m},\;
                     \underbrace{A^{(2)}_{i,j},\dots,A^{(2)}_{i,j}}_{T_2}\big).
\end{equation}

\subsection{Loss}

Plug $A_{i,j}$ into \eqref{eq:loss}, multiply the tool mask into the
model-token mask. Standard clip ($\epsilon_{\text{lo}}=0.2$, asymmetric
upper clip), $\beta=0$, token-level importance sampling. No critic, no
GAE, no context truncation across turns.

\section{Per-turn loss decomposition}
\label{sec:two-pass}

A single forward over the full trajectory
$[s\,|\,y^{(1)}\,|\,m\,|\,y^{(2)}]$ materializes an LM-head logits
gradient of shape $(B,\,T_{\text{full}},\,|V|)$. With Qwen2.5's
$|V|\approx 152{,}000$ this gradient is several gigabytes and OOMs a
24\,GB L4 at our trajectory lengths. We exploit additivity of
\eqref{eq:loss} to split the loss into two passes that share parameter
gradients but never share activations.

\paragraph{Decomposition.}
Let $\mathcal{I}_1, \mathcal{I}_2$ index turn-1 and turn-2 model
tokens, and $\mathcal{I}_m$ the tool-message tokens. Since the tool
mask zeroes the loss on $\mathcal{I}_m$ and \eqref{eq:loss} sums over
positions with a single normalizer,
\begin{equation}
  \mathcal{L} \;=\;
    \mathcal{L}_{\mathcal{I}_1} \;+\;
    \mathcal{L}_{\mathcal{I}_2},\qquad
  N_{\text{tok}} = |\mathcal{I}_1| + |\mathcal{I}_2|\ \text{shared.}
\label{eq:additive}
\end{equation}

\paragraph{Pass 1.} Forward $[s\,|\,y^{(1)}]$ with
\texttt{logits\_to\_keep}$=T_1^{\max}$, advantages $A^{(1)}_i$ over
$\mathcal{I}_1$. Compute $\mathcal{L}_1$, backward (scaled by
$1/\text{accum\_steps}$), free.

\paragraph{Pass 2.} Forward
$[s\,|\,y^{(1)}\,|\,m\,|\,y^{(2)}]$ but pack
$[s\,|\,y^{(1)}\,|\,m]$ into the prompt slot and $y^{(2)}$ alone into
the completion slot, so the trainer's
\texttt{logits\_to\_keep}$=T_2^{\max}$ and the LM-head gradient
materializes only over turn-2. Use $A^{(2)}_{i,j}$, compute
$\mathcal{L}_2$, backward, free.

By \eqref{eq:additive} the accumulated parameter gradients sum to the
single-pass gradient exactly. Peak memory is
$\max(\text{pass 1, pass 2})$, since pass-1's graph is freed before
pass-2's forward begins.

\paragraph{Restrictions.}
Equation~\eqref{eq:additive} requires the token-summed,
single-normalizer form of \eqref{eq:loss}; per-row averaging variants
do not factor across $\mathcal{I}_1\cup\mathcal{I}_2$. The
implementation also needs $\beta=0$ (a KL term to $\pi_{\text{ref}}$
would require per-segment slicing of reference log-probabilities) and
\texttt{num\_iterations}$=1$ (so cached
\texttt{old\_per\_token\_logps} need no re-segmenting across PPO-style
iterations).

\section{Related work}

The closest existing algorithm is
\emph{VinePPO}~\cite{kazemnejad2024vineppo}, which replaces PPO's
learned critic with on-policy MC values obtained by re-branching from
chosen intermediate states inside a single chain-of-thought. The
granularity is the difference: VinePPO branches at arbitrary
intra-completion token positions and uses $\widehat V(s_t)$ as a
\emph{per-token} baseline that varies within a completion. We branch
only at \emph{turn boundaries} and broadcast the resulting advantage
as a constant across all tokens of a turn. Within a turn the update is
group-relative in the GRPO sense; across turns it is VinePPO-like.
Branching depth is fixed at one ($G_1\times G_2$ tree per prompt),
matching the chat-template tool-calling protocol the policy uses at
inference time.

A second neighbor is \emph{ArCHer}~\cite{zhou2024archer}, which uses
the same macro-step decomposition (per-turn value, per-token policy
update with the turn-level advantage broadcast across action tokens)
but estimates the per-turn value with a learned critic. We estimate
it by MC over $G_2$ continuations --- the standard variance/bias
trade-off of GAE-style versus REINFORCE-style estimates. We take the
MC route because $N=2$ keeps MC variance manageable and avoids
fitting a value head alongside the LoRA-adapted policy. A related
no-critic alternative is ReBN~\cite{liu2025gem}, which replaces GAE
with batch normalization of MC returns; the hierarchical mean we use
can be viewed as ReBN with the batch partitioned by turn-1 prefix.

The same $G_1\times G_2$ tree appears in process-reward-model
training pipelines such as
Math-Shepherd~\cite{wang2024mathshepherd} and
OmegaPRM~\cite{luo2024omegaprm}: there the per-prefix MC mean labels
intermediate steps for a learned reward model. We use it as a
policy-gradient baseline rather than as a labelling oracle.

In the broader multi-turn LLM-RL literature, RAGEN's
StarPO~\cite{wang2025ragen} is the canonical critic-free MC baseline:
$N$ trajectories per prompt, flat group baseline over their rewards,
single advantage broadcast across all turns. Setting $G_2=1$
collapses our scheme to a flat per-prompt baseline of this form.
Critic-based turn-level training schemes~\cite{chen2025turn} take the
opposite route, fitting GAE on top of a learned value head.

The artifact we train is a \emph{code policy} ---
\texttt{def policy(state)} executed without further LLM calls in the
inner control loop. This matches systems like
Voyager~\cite{wang2024voyager} (frozen-LLM prompting + skill-library
accumulation) and FunSearch~\cite{romera-paredes2024funsearch} (LLM
as evolutionary mutation operator). Unlike either, we adapt the LLM
by gradient-based RL of LoRA weights and use a scratchpad turn
between prompt and committed policy.

To our knowledge, the specific combination ``turn-boundary branching
+ critic-free MC baseline + within-turn token broadcast'' is not
written down in prior LLM-RL work in this exact form, but each
ingredient is standard and a reader familiar with VinePPO and
ArCHer would not find the specialization surprising. We present it
as a clean specialization, not a new algorithm.

\section{Summary of contributions}

We do not propose a new RL algorithm. The training scheme is
VinePPO-style branched MC value estimation with a token-summed
clipped surrogate, restricted to $N=2$ chat-template turns and
applied to code-policy generation. Engineering contributions:
\begin{enumerate}
  \item the per-turn loss decomposition of \S\ref{sec:two-pass},
        which exploits additivity of \eqref{eq:loss} to bound peak
        LM-head logits memory to $\max(\text{turn 1, turn 2})$
        instead of their sum, enabling $N=2$ branched rollouts on a
        24\,GB GPU;
  \item the typed-fence two-turn protocol with a \texttt{main}
        scratchpad turn whose stdout is fed back as a tool message
        before the final \texttt{policy} commit.
\end{enumerate}

\begin{thebibliography}{99}
\bibitem{shao2024deepseekmath} Z.~Shao et~al.
\emph{DeepSeekMath: Pushing the Limits of Mathematical Reasoning in
Open Language Models}. arXiv:2402.03300, 2024.
\bibitem{yu2025dapo} Q.~Yu et~al.
\emph{DAPO: An Open-Source LLM Reinforcement Learning System at Scale}.
arXiv:2503.14476, 2025.
\bibitem{kazemnejad2024vineppo} A.~Kazemnejad et~al.
\emph{VinePPO: Unlocking RL Potential For LLM Reasoning Through Refined
Credit Assignment}. arXiv:2410.01679, 2024.
\bibitem{wang2025ragen} Z.~Wang et~al. \emph{RAGEN: Training Agents by
Reinforcing Reasoning}. 2025.
\bibitem{chen2025turn} W.-T.~Chen et~al. \emph{Turn-level training for
multi-turn LLM agents}. 2025.
\bibitem{liu2025gem} Z.~Liu et~al. \emph{GEM: Reinforcement Learning
for Generalist Models}. 2025.
\bibitem{zhou2024archer} Y.~Zhou et~al. \emph{ArCHer: Training
Language Model Agents via Hierarchical Multi-Turn RL}. ICML, 2024.
\bibitem{wang2024mathshepherd} P.~Wang et~al. \emph{Math-Shepherd:
Verify and Reinforce LLMs Step-by-step without Human Annotations}.
ACL, 2024.
\bibitem{luo2024omegaprm} L.~Luo et~al. \emph{Improve Mathematical
Reasoning in Language Models by Automated Process Supervision}.
arXiv:2406.06592, 2024.
\bibitem{wang2024voyager} G.~Wang et~al.
\emph{Voyager: An Open-Ended Embodied Agent with Large Language Models}.
TMLR, 2024.
\bibitem{romera-paredes2024funsearch} B.~Romera-Paredes et~al.
\emph{Mathematical discoveries from program search with large language
models}. Nature, 2024.
\end{thebibliography}

\end{document}
